\documentclass[11pt]{article}
\usepackage[T1]{fontenc}
\usepackage[margin=0.9in]{geometry}
\usepackage{amsmath,amssymb,booktabs,array}
\usepackage[hidelinks]{hyperref}
\newcommand{\Normal}{\mathcal N}
\newcommand{\BF}{\mathrm{BF}}
\newcommand{\E}{\mathbb E}
\newcommand{\KL}{\operatorname{KL}}
\newcommand{\ERSS}{\operatorname{ERSS}}
\newcommand{\ind}{\mathbf 1}
\newcommand{\T}{^{\mathsf T}}
\newcommand{\yb}{\widetilde{\mathbf y}}
\newcommand{\etal}{\overline{\boldsymbol\eta}}
\setlength{\parskip}{0.35em}
\setlength{\parindent}{0pt}
\allowdisplaybreaks
\title{SuSiE with a finite slider prior\\\large A 17-point model and variational EM algorithm}
\author{Model draft for \texttt{susie\_slide\_prior}}
\date{27 September 2026}
\begin{document}
\maketitle

\section{Model and the parameter learned by EM}
We extend the slider model in Denault's manuscript~\cite{slider},
Sections 2--5: retain the Gaussian prior on the effect size $\beta$,
put a discrete prior on the slider $\delta$, and learn its probabilities
by empirical Bayes. The support stays fixed:
\begin{equation}
 K=17,\qquad d_k=-1+\frac{k}{8},\quad k=0,\ldots,16.
 \label{eq:grid}
\end{equation}
Thus the spacing is $1/8=0.125$, and
\begin{align*}
\mathcal D=\{&-1,-0.875,-0.75,-0.625,-0.5,-0.375,-0.25,-0.125,0,\\
             &0.125,0.25,0.375,0.5,0.625,0.75,0.875,1\}.
\end{align*}
Write $\Pr(\delta=d_k)=w_k$, with $w_k\geq0$ and
$\sum_{k=0}^{16}w_k=1$. Initialize $w_k^{(0)}=1/17$.
Equal probabilities are the initialization; if they are fixed permanently,
there is no slider-prior parameter for EM to estimate. This is a discrete
distribution, not quadrature for a continuous uniform prior. In particular,
the endpoints initially receive the same mass as every other point.

For hard-call genotypes $x_{ij}\in\{0,1,2\}$, set
$h_{ij}=\ind\{x_{ij}=1\}$. The coding is
$f_\delta(x)=x+\delta\ind\{x=1\}$: $\delta=-1,0,1$ correspond to
recessive, additive and dominant effects relative to the counted allele.
Use the manuscript's intercept-integrated coordinates:
\begin{equation}
 \yb=U\T\mathbf y,\qquad
 \mathbf a_j=\frac{U\T\mathbf x_j}{s_j},\qquad
 \mathbf h_j=\frac{U\T\mathbf h_j^{\rm raw}}{s_j},\qquad
 \mathbf z_{jk}=\mathbf a_j+d_k\mathbf h_j.
 \label{eq:design}
\end{equation}
Here $U\T U=I_N$, $UU\T=I_n-\mathbf1\mathbf1\T/n$ and $N=n-1$.
An explicit $U$ is unnecessary: centered inner products give the same
results. Without an intercept, take $U=I_n$ and $N=n$.
Set $s_j=1$ for raw units, or use the original additive genotype standard
deviation, as in \texttt{susieSlide}. Exclude constant columns in the
latter case. Keep $s_j$ fixed across all slider values.

For $\ell=1,\ldots,L$, independently across components, let
\begin{align}
 J_\ell&\sim\operatorname{Categorical}(\pi_1,\ldots,\pi_p),&
 C_\ell&\sim\operatorname{Categorical}(w_0,\ldots,w_{16}),\nonumber\\
 \delta_\ell&=d_{C_\ell},&\beta_\ell&\sim\Normal(0,V_\ell),\nonumber\\
 \boldsymbol\eta_\ell&=\beta_\ell\mathbf z_{J_\ell C_\ell},&
 \yb&=\sum_{\ell=1}^L\boldsymbol\eta_\ell+
       \boldsymbol\epsilon,\quad
       \boldsymbol\epsilon\sim\Normal(0,\sigma^2 I_N).
 \label{eq:model}
\end{align}
The three latent variables in a component are independent a priori.
The SNP weights $\pi_j$ are fixed (usually $1/p$); $V_\ell>0$ and
$\sigma^2>0$. A single common $\mathbf w$ is learned across components.
Each component selects one SNP and one slider; there are not $p$
independent observed slider draws per component. Different components may
select the same SNP with different sliders.

\newpage
\section{Exact single-effect posterior on the grid}
For an expected residual vector $\mathbf r_\ell$, precompute
\begin{align}
 A_j&=\mathbf a_j\T\mathbf a_j,&
 B_j&=\mathbf a_j\T\mathbf h_j,&
 D_j&=\mathbf h_j\T\mathbf h_j,\nonumber\\
 t_{\ell j}&=\mathbf a_j\T\mathbf r_\ell,&
 u_{\ell j}&=\mathbf h_j\T\mathbf r_\ell.
 \label{eq:stats}
\end{align}
The first three quantities depend only on the genotypes. For all 17 values,
\begin{equation}
 S_{jk}=A_j+2d_kB_j+d_k^2D_j,\qquad
 T_{\ell jk}=t_{\ell j}+d_ku_{\ell j}.
 \label{eq:ST}
\end{equation}
Consequently, each grid evaluation needs only scalar arithmetic after two
residual inner products per SNP. No transformed $n\times17p$ matrix is needed.

Integrating the Gaussian effect size gives the conditional Bayes factor
against a zero-effect component:
\begin{align}
 v_{\ell jk}&=\left(V_\ell^{-1}+S_{jk}/\sigma^2\right)^{-1},&
 m_{\ell jk}&=v_{\ell jk}T_{\ell jk}/\sigma^2,\label{eq:moments}\\
 b_{\ell jk}:=\log\BF_{\ell jk}
 &=-\frac12\log\left(1+\frac{V_\ell S_{jk}}{\sigma^2}\right)
 +\frac{V_\ell T_{\ell jk}^{\,2}}
 {2\sigma^2(\sigma^2+V_\ell S_{jk})}.
 \label{eq:bf}
\end{align}
These are the manuscript's Gaussian formulas evaluated at $\delta=d_k$.
For one SER problem the evidence, relative to its null density $p_0$, is
\begin{equation}
 Z_\ell=\sum_{j=1}^p\sum_{k=0}^{16}\pi_jw_k\BF_{\ell jk},
 \qquad p(\mathbf r_\ell)=p_0(\mathbf r_\ell)Z_\ell.
 \label{eq:evidence}
\end{equation}
The exact SER posterior is
\begin{align}
 \rho_{\ell jk}
 &:=\Pr(J_\ell=j,C_\ell=k\mid\mathbf r_\ell)
   =\frac{\pi_jw_k\exp(b_{\ell jk})}{Z_\ell},\label{eq:rho}\\
 \beta_\ell\mid J_\ell=j,C_\ell=k,\mathbf r_\ell
 &\sim\Normal(m_{\ell jk},v_{\ell jk}).\label{eq:beta}
\end{align}
Normalize $\log\pi_j+\log w_k+b_{\ell jk}$ over \emph{all} $17p$ pairs
using log-sum-exp. The SNP evidence is the prior average
$\sum_k w_k\BF_{\ell jk}$, not the largest grid Bayes factor.

The posterior SNP probability and conditional slider probabilities are
\begin{equation}
 \alpha_{\ell j}=\sum_k\rho_{\ell jk},\qquad
 \Pr(\delta_\ell=d_k\mid J_\ell=j,\mathbf r_\ell)
       =\rho_{\ell jk}/\alpha_{\ell j},
 \label{eq:marginals}
\end{equation}
when $\alpha_{\ell j}>0$. Marginalizing $\beta$ over the slider therefore
produces a mixture of 17 Gaussians, not generally a single Gaussian.

For efficient fitted values define the \emph{joint} first moments
\begin{equation}
 a_{\ell j}^{*}=\sum_k\rho_{\ell jk}m_{\ell jk},\qquad
 h_{\ell j}^{*}=\sum_k\rho_{\ell jk}d_km_{\ell jk},\qquad
 \etal_\ell=X_c\mathbf a_\ell^{*}+H_c\mathbf h_\ell^{*},
 \label{eq:fit}
\end{equation}
where $X_c$ and $H_c$ have columns $\mathbf a_j$ and $\mathbf h_j$.
In general, $\E(\beta_\ell\delta_\ell\mid J_\ell=j)$ is not
$\E(\beta_\ell\mid J_\ell=j)\E(\delta_\ell\mid J_\ell=j)$.
Using the latter product would give an incorrect component fit.

\newpage
\section{SuSiE fitting: variational EM with IBSS updates}
Use the factorization
\begin{equation}
 q(J,C,\beta)=\prod_{\ell=1}^L q_\ell(J_\ell,C_\ell,\beta_\ell).
 \label{eq:q}
\end{equation}
Within a component retain all dependence between SNP identity, slider and
effect size. Only the factors across components are separated. As in
SuSiE~\cite{susie}, this yields a variational approximation for $L>1$;
it is not an exact E-step for the full multi-effect posterior.

\paragraph{E-step (coordinate ascent).}
For each component, restore its old fitted contribution to the current
residual and form
\begin{equation}
 \mathbf r_\ell=\yb-\sum_{b\ne\ell}\etal_b.
 \label{eq:residual}
\end{equation}
Compute equations~\eqref{eq:stats}--\eqref{eq:rho} and update
$q_\ell$ by equation~\eqref{eq:beta}. Then refresh its fitted contribution
using equation~\eqref{eq:fit} before moving to the next component.

\paragraph{M-step (learn the slider probabilities).}
With the variational distribution fixed, the part of the expected
complete-data log likelihood depending on $\mathbf w$ is
\begin{equation}
 Q(\mathbf w)=\sum_{k=0}^{16}N_k\log w_k+\text{constant},
 \qquad N_k=\sum_{\ell=1}^L\sum_{j=1}^p\rho_{\ell jk}.
 \label{eq:Q}
\end{equation}
Since $\sum_kN_k=L$, maximizing on the probability simplex gives
\begin{equation}
 \boxed{\displaystyle
   w_k^{\rm new}=\frac1L\sum_{\ell=1}^L\sum_{j=1}^p\rho_{\ell jk},
   \qquad k=0,\ldots,16.}
 \label{eq:wupdate}
\end{equation}
The denominator is $L$, not $Lp$ or $17L$. The grid values $d_k$ do not
move. Share $\mathbf w$ across components; independently fitting one prior
per candidate SNP would defeat the intended pooling.

\paragraph{Algorithm.}
\begin{enumerate}
 \item Initialize $w_k=1/17$, positive variances and zero component
 fitted means. Equivalently initialize each variational factor to its
 prior. Precompute $A_j,B_j,D_j$.
 \item Holding $\mathbf w$ and the variances fixed, cycle over
 $\ell=1,\ldots,L$ using the residual SER update above. One sweep is
 a valid coordinate-ascent pass; more sweeps may improve the E-step.
 \item Update all $w_k$ once using the current factors and
 equation~\eqref{eq:wupdate}. Optionally update $V_\ell$ and $\sigma^2$
 by the formulas on the next page.
 \item Repeat until the ELBO and prior weights stabilize. Refresh the
 posterior factors after the final hyperparameter change so that
 reported summaries agree with the fitted parameters.
\end{enumerate}

\paragraph{Why the updates are valid.}
For fixed parameters and other factors, independence in~\eqref{eq:q}
gives
\begin{equation}
 \E_q\left\|\yb-\sum_b\boldsymbol\eta_b\right\|^2
 =\E_{q_\ell}\|\mathbf r_\ell-\boldsymbol\eta_\ell\|^2+C_{-\ell},
 \label{eq:residual-identity}
\end{equation}
where $C_{-\ell}$ is independent of $q_\ell$. Hence the ELBO as a function
of $q_\ell$ equals $C+\log p(\mathbf r_\ell)
-\KL(q_\ell\|p_{\rm SER,\ell})$. The grid SER posterior minimizes this
KL to zero. Equation~\eqref{eq:wupdate} maximizes the same ELBO over
$\mathbf w$ at fixed $q$. Exact coordinate updates are therefore
nondecreasing in the ELBO, without guaranteeing a global optimum.
For $L=1$, the posterior is exact and the procedure is ordinary EM.

\newpage
\section{Variance updates, objective and posterior summaries}
The Gaussian prior on $\beta$ remains in the model. If its variance is
learned, the EM update at fixed $q$ is
\begin{equation}
 V_\ell^{\rm new}=\sum_{j,k}\rho_{\ell jk}
                         (v_{\ell jk}+m_{\ell jk}^{,2}).
 \label{eq:Vupdate}
\end{equation}
All posterior moments on the right-hand sides of the M-step are held
fixed at their current values. Do not omit uncertainty in the slider or
in the effect size when computing the residual variance:
\begin{align}
 R&=\left\|\yb-\sum_\ell\etal_\ell\right\|^2,
 \nonumber\\
 \ERSS&=R+\sum_\ell\left\{
  \sum_{j,k}\rho_{\ell jk}(v_{\ell jk}+m_{\ell jk}^{,2})S_{jk}
                         -\|\etal_\ell\|^2\right\},\label{eq:erss}\\
 (\sigma^2)^{\rm new}&=\ERSS/N.\label{eq:sigma}
\end{align}
Here $N=n-1$ is correct for the manuscript's intercept-integrated
likelihood. Matching the current package's convention of estimating a
fixed intercept in the full likelihood instead uses $n$ in this
denominator and in the likelihood normalizing term. The centered SER
updates at fixed $\sigma^2$ coincide under these two conventions.

A computable ELBO for the stated intercept-integrated model is
\begin{align}
 \mathcal F&=-\frac N2\log(2\pi\sigma^2)-\frac{\ERSS}{2\sigma^2}
             -\sum_\ell K_\ell,\label{eq:elbo}\\
 K_\ell&=\sum_{j,k}\rho_{\ell jk}
               \log\frac{\rho_{\ell jk}}{\pi_jw_k}\nonumber\\
 &\quad+\frac12\sum_{j,k}\rho_{\ell jk}
 \left\{\frac{v_{\ell jk}+m_{\ell jk}^{,2}}{V_\ell}
                 -1+\log\frac{V_\ell}{v_{\ell jk}}\right\}.
 \label{eq:kl}
\end{align}
This includes categorical uncertainty in both SNP and slider. Terms with
zero posterior mass contribute zero; a posterior cannot assign positive
mass where the prior has zero mass. Use log probabilities numerically.

For positive-variance components, ordinary SNP PIPs remain
\begin{equation}
 \operatorname{PIP}_j=1-\prod_{\ell=1}^L(1-\alpha_{\ell j}).
 \label{eq:pip}
\end{equation}
Construct component credible sets from the SNP probabilities
$\alpha_{\ell j}$ after summing over the grid. Conditional slider
summaries include
\begin{align}
 \E(\delta_\ell\mid J_\ell=j,\mathbf y)
     &\simeq\frac{\sum_k d_k\rho_{\ell jk}}{\alpha_{\ell j}},\nonumber\\
 \Pr(\delta_\ell=0\mid J_\ell=j,\mathbf y)
     &\simeq\frac{\rho_{\ell j8}}{\alpha_{\ell j}}.
 \label{eq:delta-summary}
\end{align}
These are summaries of the component-specific slider, not a unique
SNP-wide slider if several components select the same SNP.
Inference averages over slider uncertainty at the fitted $\mathbf w$;
it does not integrate uncertainty in the empirical-Bayes estimate of
$\mathbf w$. Setting $w_8=1$ exactly recovers additive SuSiE under
the same scaling, priors and intercept convention.

\newpage
\section{Computational cost relative to Gaussian additive SuSiE}
The comparison is with individual-level additive SuSiE on the same data,
with a Gaussian prior on each selected coefficient and the same $L$.
Both models integrate $\beta$ analytically.

\paragraph{Setup and one residual SER calculation.}
Both designs can be prepared in $O(np)$ work. The slider additionally
needs the heterozygote indicators and the fixed cross-products
$B_j,D_j$; these are computed once. At each new residual the additive
calculation requires $t_{\ell j}$, while the slider also requires
$u_{\ell j}$. This is one extra length-$n$ inner product per SNP per
component update: approximately $n$ multiply-accumulate operations
($2n-1$ separate multiplication/addition operations for a dense dot
product), plus 16 extra grid-specific BF/moment evaluations and
the associated mixture sums. Binary genotype implementations may use
cheaper group sums, so this is a dense arithmetic comparison.

After the two inner products, all 17 grid values cost $O(17)$ per SNP,
not $O(17n)$. If only $\mathbf w$ is changed with a fixed residual,
$V_\ell$ and $\sigma^2$, cached BFs can be reweighted in $O(17p)$
without another pass through the samples.

\paragraph{A full component update also refreshes fitted values.}
Equation~\eqref{eq:fit} adds one multiplication by $H_c$ to the usual
multiplication by $X_c$. The main dense operations are therefore
\begin{center}
\small
\begin{tabular}{@{}p{0.29\textwidth}p{0.29\textwidth}p{0.32\textwidth}@{}}
\toprule
Work per component & Additive Gaussian SuSiE & 17-point slider SuSiE\\
\midrule
Residual scores & $X_c\T\mathbf r_\ell$ &
 $X_c\T\mathbf r_\ell$, $H_c\T\mathbf r_\ell$\\[0.4em]
Scalar posterior work & $p$ candidate calculations & $17p$ pair calculations\\[0.4em]
Fitted contribution & $X_c\mathbf a_\ell^{*}$ &
 $X_c\mathbf a_\ell^{*}+H_c\mathbf h_\ell^{*}$\\[0.4em]
Grid-prior update & None & $O(17p)$ count accumulation\\
\bottomrule
\end{tabular}
\end{center}
Counting a dense multiplication and addition separately, the leading
matrix-vector work is approximately $4np$ operations per additive
component update versus $8np$ per slider component update. Thus the
extra leading cost is approximately $4np$ per component, or $4nLp$
per sweep, plus $O(17Lp)$ scalar work. Equivalently there are two
extra length-$n$ multiply-accumulate passes per SNP per full update.
Updating $\mathbf w$ itself needs no new sample-level calculation.

For $T$ total IBSS sweeps the optimized time order is
\begin{equation}
 O\{np+TL(np+17p)\},
 \label{eq:cost}
\end{equation}
compared with $O\{np+TL(np+p)\}$ for additive SuSiE. With fixed $17$
and $n$ large, both are $O(TLnp)$; the slider roughly doubles the
dominant matrix work. This is not a measured twofold runtime guarantee:
transcendental functions, implementation overhead, storage, genotype
sparsity and the number of convergence sweeps also matter.

Compared with the existing plug-in slider code, both residual products
and both fitted-value products are already present. The main difference
is replacing the cubic maximization by 17 BF/moment calculations and
mixture bookkeeping, with no additional asymptotic sample pass.
The manuscript's reported 10--20-fold integration slowdown describes
its earlier implementation; it is not an inherent cost of this finite
grid algorithm. Benchmarking is needed to compare elapsed times.

\newpage
\section{Choices relevant to implementation in this branch}
\paragraph{Learning a distribution needs replicated effects.}
There are 16 free prior weights. A single region with a small number of
signals provides limited information about their population distribution.
In particular, for $L=1$ and fixed variances,
\begin{equation}
 Z(\mathbf w)=\sum_k w_k B_k,\qquad B_k=\sum_j\pi_j\BF_{jk}.
\end{equation}
Its maximum over the simplex puts all mass on a maximizing $k$
(or a mixture of tied maxima). Thus unregularized EB with one effect
does not generally learn a diffuse distribution. To learn a shared
inheritance-mode distribution, pool independent loci $g$:
\begin{equation}
 w_k^{\rm new}=
 \frac{\sum_g\sum_{\ell=1}^{L_g}\sum_j\rho_{g\ell jk}}
      {\sum_g L_g}.
 \label{eq:pooled}
\end{equation}
Each locus retains its own residuals, noise variance and effect variances.
An optional regularized alternative is a symmetric Dirichlet prior
with parameters $1+\lambda/17$, giving the MAP update
$(N_k+\lambda/17)/(L+\lambda)$. This changes the objective to
penalized EB and is not required for the unregularized model above.

\paragraph{The current \texttt{min\_obs} rule.}
The package forces $\delta=0$ when any genotype class is too small.
If retaining this rule, let $\mathcal J_{\rm free}$ be the SNPs passing
the genotype-count filter. Use the conditional prior
\begin{equation}
 w_{jk}=\begin{cases}w_k,&j\in\mathcal J_{\rm free},\\
                       \ind\{k=8\},&j\notin\mathcal J_{\rm free}.
          \end{cases}
 \label{eq:forced-prior}
\end{equation}
Replace $w_k$ by $w_{jk}$ in the SER probabilities and categorical KL.
A forced SNP has evidence $\BF_{\ell j8}$, without a factor $w_8$.
It does not provide evidence that the learned population prior is
additive. The correct update becomes
\begin{equation}
 w_k^{\rm new}=
 \frac{\sum_\ell\sum_{j\in\mathcal J_{\rm free}}\rho_{\ell jk}}
      {\sum_\ell\sum_{j\in\mathcal J_{\rm free}}\alpha_{\ell j}}.
 \label{eq:forced-update}
\end{equation}
If the denominator is zero, retain the previous weights. The same
restriction applies to pooled-locus sums. Disabling the count rule
gives the simpler common-prior model and update~\eqref{eq:wupdate}.

\paragraph{Zero effects, orientation and storage.}
The main derivation assumes $V_\ell>0$. An exactly zero-variance
component can instead be represented as the deterministic zero vector
with no SNP or slider draw, and omitted from prior-learning counts and
the count denominator. A null alternative similarly has no slider
draw. Weak positive-variance components remain in the stated model;
dropping them by a heuristic threshold changes the update.

Reversing the counted allele maps $(\beta,\delta)$ to
$(-\beta,-\delta)$, after changing the intercept. A learned asymmetric
prior therefore needs a consistent allele convention. If invariance
to arbitrary per-SNP allele flips is desired, impose $w_k=w_{16-k}$;
the constrained M-step pools each reflected pair's counts equally.

Storing every posterior pair takes $O(17Lp)$ memory. This can be
reduced by processing a component at a time and retaining its
aggregated moments and grid counts; full pair posteriors can be
recomputed when requested. $H_c$ can also be formed in chunks.
Neither an $n\times17p$ design nor 17 residual cross-products is needed.

\begin{thebibliography}{9}
\bibitem{slider} W. Denault (September 2026).
\emph{The slider model}. User-supplied manuscript
\texttt{slider\_model-6.pdf}, especially Sections 2--5 and 8.
\bibitem{susie} G. Wang, A. Sarkar, P. Carbonetto and M. Stephens (2020).
A simple new approach to variable selection in regression, with
application to genetic fine mapping. \emph{JRSS B}, 82, 1273--1300.
\href{https://doi.org/10.1111/rssb.12388}{doi:10.1111/rssb.12388}.
\end{thebibliography}
\end{document}
